\documentclass[11pt,a4paper]{letter}
\usepackage[utf8]{inputenc}
\usepackage[T1]{fontenc}
\usepackage{microtype}
\usepackage[margin=1in]{geometry}
\usepackage{hyperref}

\signature{Clarimar Jos\'{e} Coelho \\
           (on behalf of all authors) \\
           \texttt{clarimar@pucgoias.edu.br}}

\address{Clarimar Jos\'{e} Coelho \\
         Pontif\'{i}cia Universidade Cat\'{o}lica de Goi\'{a}s \\
         Av. Universit\'{a}ria, 1069 --- Setor Universit\'{a}rio \\
         Goi\^{a}nia, GO 74605--010, Brazil}

\begin{document}

\begin{letter}{Professor Federico Marini \\
               Editor-in-Chief \\
               \textit{Chemometrics and Intelligent Laboratory Systems} \\
               Elsevier}

\opening{Dear Professor Marini,}

We are pleased to submit the revised version of manuscript
\textbf{CHEMOLAB-D-26-00658R1}, entitled
\textit{``Prior-Guided Spectral Gating for Interpretable Deep Learning in
Near-Infrared Spectroscopy,''} for your consideration.

The revision addresses the single point raised by Reviewer~\#2 in the minor
revision decision (Comment~1, Experiments), concerning the risk of
overfitting in the selection of the number of PLS components. A detailed
point-by-point response to the reviewer is provided separately in the file
\texttt{response\_letter\_R2}. Below we summarise the modifications made
to the manuscript itself so that the editorial office can quickly locate
each change.

\medskip
\textbf{Summary of the technical change.}
We now adopt the \emph{one-standard-error rule} of Hastie, Tibshirani and
Friedman (2009) as a pre-registered criterion for selecting the number of
PLS latent variables~$H$, applied to the leave-one-out cross-validation
curve on the training set. This is the standard formal safeguard against
the overfitting risk correctly identified by the reviewer.

Applied \emph{ex post} to the identical training partitions used in the
original submission (verified by exact reproduction of the R1 PLS
baselines: Tecator $R^2=0.9191$, Gasoline $R^2=0.9414$), the one-SE rule
yields:

\begin{itemize}
  \item \textbf{Tecator} ($n_{\text{train}}=172$): the argmin of RMSECV
        lies at $H=14$, but the one-SE rule selects $H^*=10$---exactly
        the value originally reported. All Tecator numbers in the
        manuscript remain unchanged.
  \item \textbf{Gasoline} ($n_{\text{train}}=48$): the argmin lies at
        $H=6$, and the one-SE rule selects $H^*=5$ (not $H=10$ as
        originally reported). Applied to the identical train/test split,
        PLS Gasoline improves from $R^2=0.941$ (RMSE\,=\,0.223) to
        $R^2=0.975$ (RMSE\,=\,0.145).
\end{itemize}

None of the four qualitative conclusions of the manuscript is affected:
(i)~the Plain CNN failure is unchanged (it does not depend on the PLS
baseline); (ii)~the identification of 931~nm as the dominant learned
C--H feature is unchanged (it derives from our model's gate weights, not
from PLS); (iii)~the lightweight-architecture argument is unchanged;
(iv)~the $R^2>0.85$ threshold discussion remains valid for both
baselines. The predictive gap on Gasoline nevertheless widens from
$1.3\%$ to $4.7\%$, and three paragraphs were reworded accordingly to
reflect the fair comparison against a stronger PLS baseline.

\medskip
\textbf{List of modifications to the manuscript.}
Nine cirurgical edits were applied, in the order they appear:

\begin{enumerate}
  \setlength\itemsep{0.1em}
  \item \textbf{Abstract} (second sentence): Gasoline PLS comparator
        updated from ``$R^2=0.93$ vs $0.94$'' to
        ``$R^2=0.93$ vs $0.98$''.
  \item \textbf{Section~5.2 (Evaluation Protocol)}: the paragraph on
        PLS component selection has been rewritten to describe the
        one-standard-error rule, with citation to
        Hastie~et~al.~(2009), and reports the RMSECV values, standard
        errors, and $H^*$ selections for both datasets.
  \item \textbf{Section~4.3 (Baselines)}: the phrase ``PLS with 10
        components, selected via cross-validation'' has been replaced by
        ``PLS with $H^*$ components selected by the one-standard-error
        rule of Hastie~et~al.\ ($H^*=10$ for Tecator, $H^*=5$ for
        Gasoline).''
  \item \textbf{Table~1}: the Gasoline PLS row has been updated from
        RMSE\,=\,0.223, $R^2=0.941$ (H\,=\,10) to
        RMSE\,=\,0.145, $R^2=0.975$ (H\,=\,5). All other rows unchanged.
  \item \textbf{Figure~2 caption}: Gasoline comparator updated
        ($R^2=0.928$ vs $0.975$; gap $4.7\%$).
  \item \textbf{Section~6.2 (Multi-Dataset Performance)}: the paragraph
        discussing Gasoline results has been reworded to reflect the
        widened gap and to clarify that the added complexity of the
        proposed framework is justified by interpretability rather than
        predictive superiority.
  \item \textbf{Figure~3 caption}: the ``$86$--$93\%$ of PLS
        performance'' descriptor has been updated to
        ``$86$--$95\%$''.
  \item \textbf{Section~7 (Addressing Potential Criticisms)}: the third
        bullet has been reworded; the argument ``in Gasoline the gap is
        only $1.3\%$, suggesting that in some domains our framework
        approaches parity with the chemometric gold standard'' has been
        replaced by an honest acknowledgment that the gap widens to
        $4.7\%$ under fair comparison against a formally regularized
        baseline.
  \item \textbf{Section~8 (Conclusion), first paragraph}: PLS
        comparator range updated from ``$R^2=0.86$--$0.93$ vs
        $0.92$--$0.94$'' to ``$R^2=0.86$--$0.93$ vs $0.92$--$0.98$''.
\end{enumerate}

\medskip
\textbf{New bibliographic entry.}
One new reference has been added to support the one-standard-error rule:

\begin{quote}
Hastie, T., Tibshirani, R., \& Friedman, J. (2009).
\textit{The Elements of Statistical Learning: Data Mining, Inference,
and Prediction} (2nd~ed.). Springer, New York.
\end{quote}

The BibTeX entry is included in the submission package as
\texttt{hastie2009\_bib\_entry.bib}. The citation key is
\texttt{hastie2009elements}.

\medskip
\textbf{Reproducibility.}
The one-SE analysis is fully reproducible from the code deposited in
the project repository at
\url{https://github.com/clarimar/prior-guided-spectral-gating}
(folder \texttt{revision\_r1/}), which contains the analysis script
\texttt{one\_se\_analysis.py} and the corresponding output. The script
consumes the same pickles generated by
\texttt{src/data/preprocessing\_multi.py} that produced the original
Table~1 numbers, and can be re-executed by any reader wishing to verify
the $H^*$ selections in under one minute.

\medskip
\textbf{Figures.}
Figures~2 and~3 have been regenerated to reflect the new PLS Gasoline
values. All other figures are unchanged.

\medskip
We thank the reviewer for the constructive comment, which has
strengthened the methodological rigour of the manuscript by replacing
an implicit stability argument with a formally justified selection rule.

\closing{Yours sincerely,}

\end{letter}
\end{document}
